%-----------------------------------------------------------------------------
% Beginning of preface.tex
%-----------------------------------------------------------------------------
%
% AMS-LaTeX 1.2 sample file for a monograph, based on amsbook.cls.
% This is a data file input by chapter.tex.
%%%%%%%%%%%%%%%%%%%%%%%%%%%%%%%%%%%%%%%%%%%%%%%%%%%%%%%%%%%%%%%%%%%%%%%%

\chapter*{Preface}


This is a graduate-level textbook about algorithms for computing with
modular forms.  It is nontraditional in that the primary focus is not
on underlying theory; instead, it answers the question {\em ``how do you use a
computer to explicitly compute spaces of modular forms?''}

This book emerged from notes for a course the author
taught at Harvard University in 2004, a course
at UC San Diego in 2005, and a course at
the University of Washington in 2006. 

The author has spent years trying to find good practical ways to
compute with classical modular forms for congruence subgroups of
$\SL_2(\Z)$ and has implemented most of these algorithms several
times, first in C++ \cite{stein:hecke}, then in MAGMA \cite{magma},
and as part of the free open source computer
algebra system \sage{} (see \cite{sage}).  Much of this work has involved
turning formulas and constructions buried in obscure research papers
into precise computational recipes then testing these 
and eliminating inaccuracies.
%The goal of this book is to explain some of what I
%have learned along the way.
%, and also describe unsolved problems whose
%solution would move the theory forward.

The author is aware of no other textbooks on computing with modular
forms, the closest work being Cremona's book \cite{cremona:algs},
which is about computing with elliptic curves, and Cohen's book
\cite{cohen:course_ant} about algebraic number theory.  
%There are
%several Ph.D. thesis on computing modular forms, including
%\cite{}\edit{me, gabor, lassine, belgium mod p guy, maybe denis
%  charles???}.  
%Though many of the techniques described in this book
%have been around for decades, the field is still not yet mature, and
%there are missing details and potential improvements to many of the
%algorithms, which you the reader might fill in, and which would be
%greatly appreciated by other mathematicians.

In this book we focus on how to compute {\em in practice} the spaces
$M_k(N,\eps)$ of modular forms, where $k\geq 2$ is an integer
and~$\eps$ is a Dirichlet character of modulus~$N$ (the appendix
treats modular forms for higher rank groups).  We spend the most
effort explaining the general algorithms that appear so far to be the best (in
practice!) for such computations.  We will not discuss in any detail
computing with quaternion algebras,
half-integral weight forms, weight 1 forms, forms for
noncongruence subgroups or groups other than $\GL_2$, Hilbert and
Siegel modular forms, trace formulas, $p$-adic modular forms, and
modular abelian varieties, all of which are topics for additional
books.  We also rarely analyze the complexity of the
algorithms, but instead settle for occasional remarks about their
practical efficiency.  

% \edit{Actually, I probably
%will add some less in-depth discussion of each of these topics.}

%Classical modular forms, which are the main focus of this book, are
%more accessible than more general automorphic forms (see the
%appendix).  
For most of this book we assume the reader has some prior exposure to
modular forms (e.g., \cite{diamond-shurman}), though we recall many of
the basic definitions.  We cite standard books for proofs of the
fundamental results about modular forms that we will use.  The reader
should also be familiar with basic algebraic number theory, linear
algebra, complex analysis (at the level of \cite{ahlfors}), and
algorithms (e.g., know what an algorithm is and what big oh notation
means). In some of the examples and applications we assume that the
reader knows about elliptic curves at the level
of~\cite{silverman:aec}.


Chapter~\ref{ch:modform} is foundational for the rest of this book.  It
introduces congruence subgroups of $\SL_2(\Z)$ and modular forms as
functions on the complex upper half plane.  We discuss $q$-expansions,
which provide an important computational handle on modular forms.  We
also study an algorithm for computing with congruence subgroups.  The
chapter ends with a list of applications of modular forms throughout
mathematics.

In Chapter~\ref{ch:level1} we discuss level 1 modular forms in much
more detail.  In particular, we introduce Eisenstein series and the
cusp form $\Delta$ and describe their $q$-expansions and basic
properties.  Then we prove a structure theorem for level 1 modular
forms and use it to deduce dimension formulas and give an algorithm
for explicitly computing a basis.  We next introduce Hecke operators
on level 1 modular forms, prove several results about them, and deduce
multiplicativity of the Ramanujan $\tau$ function as an application.
We also discuss explicit computation of Hecke operators.  
In Section~\ref{sec:fasttau} we make some brief remarks on
recent work on asymptotically fast computation of
values of $\tau$. Finally, we describe computation of constant terms
of Eisenstein series using an analytic algorithm.  We generalize many
of the constructions in this chapter to higher level in subsequent
chapters.

In Chapter~\ref{ch:modform2} we turn to modular forms of higher level
but restrict for simplicity to weight $2$ since much is clearer in
this case.  (We remove the weight restriction later in
Chapter~\ref{ch:modsym}.)  We describe a geometric way of viewing
cuspidal modular forms as differentials on modular curves, which leads
to modular symbols, which are an explicit way to present a
certain homology group.  This chapter closes with methods for
explicitly computing cusp forms of weight 2 using modular symbols,
which we generalize in Chapter~\ref{ch:modformcomp}.

In Chapter~\ref{ch:dirichlet} we introduce Dirichlet characters, which
are important both in explicit construction of Eisenstein series (in 
Chapter~\ref{ch:eisen}) and
in decomposing spaces of modular forms as direct sums of simpler
spaces.  The main focus of this chapter is a detailed study of
how to explicitly represent and compute with Dirichlet characters.

Chapter~\ref{ch:eisen} is about how to explicitly construct the
Eisenstein subspace of modular forms.  First we define generalized
Bernoulli numbers attached to a Dirichlet character and an integer
then explain a new analytic algorithm for computing them (which
generalizes the algorithm in Chapter~\ref{ch:level1}).  Finally we
give without proof an explicit description of a basis of Eisenstein
series, explain how to compute it, and give some examples.

Chapter~\ref{ch:dim} records a wide range of dimension formulas for
spaces of modular forms, along with a few remarks about where they
come from and how to compute them.  

Chapter~\ref{ch:linalg} is about linear algebra over exact fields,
mainly the rational numbers.  This chapter can be read independently
of the others and does not require any background in modular forms.
Nonetheless, this chapter occupies a central position in this book,
because the algorithms in this chapter are of crucial importance to
any actual implementation of algorithms for computing with modular
forms.

Chapter~\ref{ch:modsym} is the most important chapter in this book; it
generalizes Chapter~\ref{ch:modform2} to higher weight and general
level.  The modular symbols formulation described here is central to
general algorithms for computing with modular forms.

Chapter~\ref{ch:modformcomp} applies the algorithms from
Chapter~\ref{ch:modsym} to the problem of computing with modular forms. First
we discuss decomposing spaces of modular forms using Dirichlet
characters, and then explain how to compute a basis of Hecke eigenforms
for each subspace using several approaches.  We also discuss
congruences between modular forms and bounds needed to provably
generate the Hecke algebra. 

Chapter~\ref{ch:periods} is about computing analytic invariants
of modular forms.  It discusses tricks for speeding convergence
of certain infinite series and sketches how to compute every
elliptic curve over $\Q$ with given conductor.

Chapter~\ref{ch:solutions} contains detailed solutions to most of the
exercises in this book.  (Many of these were written by students
in a course taught at the University of Washington.)

Appendix~\ref{ch:appendix} deals with computational techniques
for working with generalizations of modular forms to more general
groups than $\SL_2(\Z)$, such as $\SL_n(\Z)$ for $n\geq 3$.  Some of
this material requires more prerequisites than the rest of the book.
Nonetheless, seeing a natural generalization of the material in the
rest of this book helps to clarify the key ideas.  The
topics in the appendix are directly related to the main themes of
this book: modular symbols, Manin symbols, cohomology
of subgroups of $\SL_2(\Z)$ with various coefficients, explicit
computation of modular forms, etc.

\medskip

%begin_sage
\vspace{2ex}
\noindent{}{\bf Software.} 
We use \sage{}, Software for Algebra and Geometry Experimentation 
(see \cite{sage}), to illustrate how to do
many of the examples.  \sage{} is completely free and packages
together a wide range of open source mathematics software for doing
much more than just computing with modular forms.  \sage{} can be
downloaded and run on your computer or can be used via a web browser over
the Internet.  The reader is encouraged to experiment with many of the
objects in this book using \sage.  We do not describe the basics of
using \sage{} in this book; the reader should read the \sage{}
tutorial (and other documentation) available at the \sage{} website
\cite{sage}.   All examples in this book have been automatically
tested and should work exactly as indicated in \sage{} version
at least 1.5. 

%end_sage

%The text of this book is licensed under the GNU Free Documentation
%License, Version 1.2, November 2002.  This means that you may freely
%copy this book.  For the precise details, see the Appendix.

%\noindent{}{\bf Acknowledgement:} None yet.
% \mbox{}\vspace{2ex} 

% \noindent{}William Stein\\
% {\tt http://modular.fas.harvard.edu}\\
% {\tt was@math.harvard.edu}

\vspace{2ex}
\noindent{}{\bf Acknowledgements.} 
David Joyner and Gabor Wiese carefully read the book and provided
a huge number of helpful comments.

John Cremona and Kevin Buzzard both made many helpful remarks that
were important in the development of the algorithms in this
book.  Much of the mathematics (and some of the writing) in
Chapter~\ref{ch:periods} is joint work with Helena Verrill.

Noam Elkies made remarks about Chapters~\ref{ch:modform} and
\ref{ch:level1}.  S\'andor Kov\'acs provided interesting comments on
Chapter~\ref{ch:modform}.  Allan Steel provided helpful feedback on
Chapter~\ref{ch:linalg}.  Jordi Quer made useful remarks
about Chapter~\ref{ch:dirichlet} and Chapter~\ref{ch:dim}.

The students in the courses that I taught on this material at Harvard,
San Diego, and Washington provided substantial feedback: in
particular, Abhinav Kumar made numerous observations about computing
widths of cusps (see Section~\ref{sec:cuspwidth}) and Thomas James
Barnet-Lamb made helpful remarks about 
how to represent Dirichlet characters.  James
Merryfield made helpful remarks about complex analytic issues and
about convergence in Stirling's formula.  Robert Bradshaw, Andrew
Crites (who wrote \exref{ch:linalg}{ex:hnf}), Michael Goff, Dustin
Moody, and Koopa Koo wrote most of the solutions included in
Chapter~\ref{ch:solutions} and found numerous typos throughout the
book.  Dustin Moody also carefully read through the book and provided
feedback.

H.~Stark suggested using Stirling's formula in Section~\ref{sec:bernnum},
and Mark Watkins and Lynn Walling made comments on
Chapter~\ref{ch:modform2}.

Justin Walker found typos in the first published version of the book.

Parts of Chapter~\ref{ch:modform} follow Serre's beautiful
introduction to modular forms \cite[Ch.~VII]{serre:arithmetic}
closely, though we adjust the notation, definitions, and order of
presentation to be consistent with the rest of this book. 

I would like to acknowledge the partial support of 
NSF Grant DMS 05-55776.
Gunnells was supported in part by NSF Grants DMS 02-45580 and
 DMS 04-01525.


\vspace{2ex}
\noindent{\bf Notation and Conventions.}
We denote canonical isomorphisms by $\isom$
and noncanonical isomorphisms by $\ncisom$. 
If $V$ is a vector space and $s$ denotes
some sort of construction involving $V$,
we let $V_s$ denote the corresponding subspace
and $V^s$ the quotient space.  E.g., if $\iota$
is an involution of $V$, then $V_+$ is $\Ker(\iota-1)$
and $V^+ = V/\Im(\iota-1)$.  If $A$ is a finite
abelian group, then $A_{\tor}$ denotes the torsion
subgroup and $A/{\tor}$ denotes the quotient $A/A_{\tor}$.
We denote right group actions using exponential
notation.  Everywhere in this book,~$N$ is a positive
integer and $k$ is an integer. 

If $N$ is an integer, a \defn{divisor} $t$ of $N$
is a {\em positive} integer such that $N/t$ is an integer.

%-----------------------------------------------------------------------------
% End of preface.tex
%-----------------------------------------------------------------------------
